% docs/week4.tex
% Week 4: the outlier-faithfulness measure, version 1, tested on the
% exploration set. The same content as docs/MEASURE.md (sections 2, 4, 5
% and 7), shortened. Needs amsmath and amssymb only. There is no
% bibliography: references are written out in the text.
%
% Use: paste into progress.tex, or put % docs/week4.tex
% Week 4: the outlier-faithfulness measure, version 1, tested on the
% exploration set. The same content as docs/MEASURE.md (sections 2, 4, 5
% and 7), shortened. Needs amsmath and amssymb only. There is no
% bibliography: references are written out in the text.
%
% Use: paste into progress.tex, or put % docs/week4.tex
% Week 4: the outlier-faithfulness measure, version 1, tested on the
% exploration set. The same content as docs/MEASURE.md (sections 2, 4, 5
% and 7), shortened. Needs amsmath and amssymb only. There is no
% bibliography: references are written out in the text.
%
% Use: paste into progress.tex, or put % docs/week4.tex
% Week 4: the outlier-faithfulness measure, version 1, tested on the
% exploration set. The same content as docs/MEASURE.md (sections 2, 4, 5
% and 7), shortened. Needs amsmath and amssymb only. There is no
% bibliography: references are written out in the text.
%
% Use: paste into progress.tex, or put \input{week4} in main.tex.
% Numbers: results/week4_measure.json and results/week4_lattice.json,
% written by src/explore_measure.py and src/lattice_check.py.

\section*{Week 4: outlier faithfulness, version 1}

Built and tested on the exploration set only: shuttle, wbc, ionosphere,
fault and magic. The repository calls these five the development set. No
benchmark dataset was used. Version 1 will change in week 5 and is frozen
in week 12.

\subsection*{Definition}

Let $x_1, \dots, x_n \in \mathbb{R}^D$ be the data and
$y_1, \dots, y_n \in \mathbb{R}^d$ a projection of it. The same detector
scores both spaces: $a_i$ is the distance from $x_i$ to its $k$-th nearest
neighbour among the other points, and $b_i$ the same for $y_i$, with
$k = 20$. Outlying in the projection therefore means ranked high by the
detector run on the projected coordinates.

\paragraph{Ranks.} $\rho_a(i)$ is the descending rank of $a_i$, with
$\rho_a(i) = 1$ for the most outlying point. Scores within a relative
tolerance of $10^{-9}$ count as tied, and tied points share the average of
the positions they span. $\rho_b$ is defined the same way.

\paragraph{Displacement,} in bits:
\begin{equation*}
  \lambda_i = \log_2 \frac{\rho_a(i)}{\rho_b(i)} .
\end{equation*}
A negative value means the projection moved the point towards the crowd, a
positive value towards the edge. $|\lambda_i|$ counts rank doublings, so a
move from rank 1 to rank 64 is $-6$ and a move from rank 1001 to rank 1064
is $-0.09$. The top of the ranking, where an analyst looks, dominates.

\paragraph{Shortlists.} With budget $m = \lceil 0.05\, n \rceil$,
$S_a = \{ i : \rho_a(i) \le m \}$ and $S_b = \{ i : \rho_b(i) \le m \}$. A
tied group shares one rank, so it is wholly in or wholly out.

\paragraph{Noise null.} For $r = 1, \dots, 100$, draw a random half $H_r$
of the points and score every point by its distance to the $(k/2)$-th
nearest neighbour within $H_r$, itself excluded. Ranking all $n$ points
gives $\rho^{(r)}$. The null for point $i$ is the set of differences
$D_i^{(rs)} = \log_2 \rho^{(r)}(i) - \log_2 \rho^{(s)}(i)$ over all pairs
$r < s$, and
\begin{equation*}
  p_i = \tfrac{1}{2} \Pr_{r<s}\bigl( |D_i^{(rs)}| \ge |\lambda_i| \bigr)
\end{equation*}
is the one-sided p-value in the direction the point moved. The null asks how
far the point's rank moves between two equally valid measurements of the
original data. It does not involve the projection, so it is computed once
per dataset.

\paragraph{Calls,} at level $\alpha = 0.05$:
\begin{center}
\begin{tabular}{lll}
\hline
category & rule & meaning \\
\hline
retained & $i \in S_a \cap S_b$ & kept on the shortlist \\
masked & $i \in S_a \setminus S_b$ and $p_i < \alpha$ & hidden, beyond noise \\
phantom & $i \in S_b \setminus S_a$ and $p_i < \alpha$ & promoted, beyond noise \\
ambiguous & crossed, $p_i \ge \alpha$ & crossed within noise \\
\hline
\end{tabular}
\end{center}

\paragraph{Aggregates.}
\begin{equation*}
  \mathrm{MR} = \frac{|M|}{|S_a|}, \qquad
  \mathrm{PR} = \frac{|P|}{|S_b|}, \qquad
  \mathrm{AO} = \frac{O - E}{\min(|S_a|, |S_b|) - E},
\end{equation*}
where $M$ and $P$ are the masked and phantom points, $O = |S_a \cap S_b|$,
and $E = |S_a|\,|S_b| / n$ is the overlap expected when the projected scores
are permuted across points. MR is the share of the detector's shortlist that
the projection hides, PR the share of the projection's shortlist that is not
outlying in the data. AO is 1 for identical shortlists and 0 at chance, the
same adjustment Campos et al.\ (2016) apply to precision at $n$.

\paragraph{Labels} are not used by the measure. A labelled anomaly that is
not outlying in the original space cannot be masked, which is the
conditioning on original-space outlyingness the proposal asks for. Labels
only say what a failure costs: a masked anomaly is a detection lost to the
projection.

\subsection*{Checks with known answers}

All 44 checks in \texttt{src/validate.py} pass. Among them: an identity
projection and a rotation move nothing; swapping the two spaces negates
$\lambda$ and swaps masked with phantom; a monotone rescoring changes
nothing; shuffled projected scores give an AO of 0 on average; comparing two
fresh re-measurements of the original data calls 2 to 3\% of crossings at
$\alpha = 5\%$; and of 20 outliers planted in the directions PCA discards,
17 to 20 are called masked over five seeds, while all 20 planted in the
plane PCA keeps are retained.

\subsection*{Results in two dimensions, top 5\%}

Masking rate MR, with the adjusted overlap AO in brackets. Means over
seeds: PCA is deterministic, the Gaussian random projection has 5
seeds, and UMAP, t-SNE and PaCMAP 3 each, at library defaults.
\begin{center}
\small
\begin{tabular}{lccccc}
\hline
 & shuttle & wbc & ionosphere & fault & magic \\
\hline
PCA & 0.57 ($+0.34$) & 0.00 ($+0.91$) & 1.00 ($-0.05$) & 0.58 ($+0.37$) & 0.39 ($+0.56$) \\
Random & 0.40 ($+0.49$) & 0.37 ($+0.58$) & 0.66 ($+0.31$) & 0.59 ($+0.36$) & 0.38 ($+0.58$) \\
UMAP & 0.91 ($+0.02$) & 0.61 ($+0.35$) & 1.00 ($-0.05$) & 0.82 ($+0.11$) & 0.85 ($+0.10$) \\
t-SNE & 0.73 ($+0.18$) & 0.81 ($+0.15$) & 0.85 ($+0.10$) & 0.70 ($+0.25$) & 0.93 ($+0.02$) \\
PaCMAP & 0.85 ($+0.07$) & 0.97 ($-0.03$) & 1.00 ($-0.05$) & 0.98 ($-0.03$) & 0.89 ($+0.06$) \\
\hline
\end{tabular}
\end{center}

\begin{itemize}
  \item UMAP, t-SNE and PaCMAP hide 61 to 100\% of the top 5\%, more than PCA
    or a random projection on four of the five datasets. PaCMAP's shortlist
    is close to chance on all five (AO at most 0.09).
  \item PCA hides from 0\% (wbc) to 100\% (ionosphere). The failures fall as
    PCA keeps more dimensions and vanish by full dimension on all five.
  \item A random projection hides less than PCA on shuttle and ionosphere,
    because PCA discards low-variance directions whatever they hold.
  \item Aggregate correlations miss the shortlist. On ionosphere PCA keeps a
    Spearman correlation of 0.56, and 0.55 for the weighted correlation of
    Schubert et al.\ (2012), while its 2D shortlist shares nothing with the
    original one.
  \item Shuttle is not a control for PCA, which hides 57\% of its top 5\%:
    its median LID is about 2 locally, but it spans six linear dimensions.
    It is also lattice data: all nine features are integers, and its 2D
    pictures under PCA, UMAP, t-SNE and PaCMAP each break into
    1,188 to 1,511 small islands of near-identical points
    (169 under a random projection), against 5 to 65 for magic.
  \item The neighbour-based methods gather isolated outliers into tight
    clumps. On wbc and fault their verdict depends on the neighbourhood size
    used to read the picture; on shuttle, ionosphere and magic it does not.
\end{itemize}

\subsection*{Open, for the supervisor}

\begin{enumerate}
  \item A resampling null per point, with permutation kept only for the
    chance level of AO, because the scDEED-style permutation null calls
    almost nothing. Is that acceptable?
  \item Budget: the top 5\% of points, which needs no labels, instead of the
    number of labelled anomalies?
  \item How to read the projection: the same $k$ as the original space (the
    proposal's rule), a profile over several $k$, or a multi-scale score?
  \item Keep shuttle as the case where local and global dimension differ?
  \item Lattice data: keep shuttle and wbc as they are and report the effect?
    The discrete fields of network data will raise the same question.
  \item Ranks, as now, or normalised scores (Kriegel et al.\ 2011) for
    version 2?
\end{enumerate}
 in main.tex.
% Numbers: results/week4_measure.json and results/week4_lattice.json,
% written by src/explore_measure.py and src/lattice_check.py.

\section*{Week 4: outlier faithfulness, version 1}

Built and tested on the exploration set only: shuttle, wbc, ionosphere,
fault and magic. The repository calls these five the development set. No
benchmark dataset was used. Version 1 will change in week 5 and is frozen
in week 12.

\subsection*{Definition}

Let $x_1, \dots, x_n \in \mathbb{R}^D$ be the data and
$y_1, \dots, y_n \in \mathbb{R}^d$ a projection of it. The same detector
scores both spaces: $a_i$ is the distance from $x_i$ to its $k$-th nearest
neighbour among the other points, and $b_i$ the same for $y_i$, with
$k = 20$. Outlying in the projection therefore means ranked high by the
detector run on the projected coordinates.

\paragraph{Ranks.} $\rho_a(i)$ is the descending rank of $a_i$, with
$\rho_a(i) = 1$ for the most outlying point. Scores within a relative
tolerance of $10^{-9}$ count as tied, and tied points share the average of
the positions they span. $\rho_b$ is defined the same way.

\paragraph{Displacement,} in bits:
\begin{equation*}
  \lambda_i = \log_2 \frac{\rho_a(i)}{\rho_b(i)} .
\end{equation*}
A negative value means the projection moved the point towards the crowd, a
positive value towards the edge. $|\lambda_i|$ counts rank doublings, so a
move from rank 1 to rank 64 is $-6$ and a move from rank 1001 to rank 1064
is $-0.09$. The top of the ranking, where an analyst looks, dominates.

\paragraph{Shortlists.} With budget $m = \lceil 0.05\, n \rceil$,
$S_a = \{ i : \rho_a(i) \le m \}$ and $S_b = \{ i : \rho_b(i) \le m \}$. A
tied group shares one rank, so it is wholly in or wholly out.

\paragraph{Noise null.} For $r = 1, \dots, 100$, draw a random half $H_r$
of the points and score every point by its distance to the $(k/2)$-th
nearest neighbour within $H_r$, itself excluded. Ranking all $n$ points
gives $\rho^{(r)}$. The null for point $i$ is the set of differences
$D_i^{(rs)} = \log_2 \rho^{(r)}(i) - \log_2 \rho^{(s)}(i)$ over all pairs
$r < s$, and
\begin{equation*}
  p_i = \tfrac{1}{2} \Pr_{r<s}\bigl( |D_i^{(rs)}| \ge |\lambda_i| \bigr)
\end{equation*}
is the one-sided p-value in the direction the point moved. The null asks how
far the point's rank moves between two equally valid measurements of the
original data. It does not involve the projection, so it is computed once
per dataset.

\paragraph{Calls,} at level $\alpha = 0.05$:
\begin{center}
\begin{tabular}{lll}
\hline
category & rule & meaning \\
\hline
retained & $i \in S_a \cap S_b$ & kept on the shortlist \\
masked & $i \in S_a \setminus S_b$ and $p_i < \alpha$ & hidden, beyond noise \\
phantom & $i \in S_b \setminus S_a$ and $p_i < \alpha$ & promoted, beyond noise \\
ambiguous & crossed, $p_i \ge \alpha$ & crossed within noise \\
\hline
\end{tabular}
\end{center}

\paragraph{Aggregates.}
\begin{equation*}
  \mathrm{MR} = \frac{|M|}{|S_a|}, \qquad
  \mathrm{PR} = \frac{|P|}{|S_b|}, \qquad
  \mathrm{AO} = \frac{O - E}{\min(|S_a|, |S_b|) - E},
\end{equation*}
where $M$ and $P$ are the masked and phantom points, $O = |S_a \cap S_b|$,
and $E = |S_a|\,|S_b| / n$ is the overlap expected when the projected scores
are permuted across points. MR is the share of the detector's shortlist that
the projection hides, PR the share of the projection's shortlist that is not
outlying in the data. AO is 1 for identical shortlists and 0 at chance, the
same adjustment Campos et al.\ (2016) apply to precision at $n$.

\paragraph{Labels} are not used by the measure. A labelled anomaly that is
not outlying in the original space cannot be masked, which is the
conditioning on original-space outlyingness the proposal asks for. Labels
only say what a failure costs: a masked anomaly is a detection lost to the
projection.

\subsection*{Checks with known answers}

All 44 checks in \texttt{src/validate.py} pass. Among them: an identity
projection and a rotation move nothing; swapping the two spaces negates
$\lambda$ and swaps masked with phantom; a monotone rescoring changes
nothing; shuffled projected scores give an AO of 0 on average; comparing two
fresh re-measurements of the original data calls 2 to 3\% of crossings at
$\alpha = 5\%$; and of 20 outliers planted in the directions PCA discards,
17 to 20 are called masked over five seeds, while all 20 planted in the
plane PCA keeps are retained.

\subsection*{Results in two dimensions, top 5\%}

Masking rate MR, with the adjusted overlap AO in brackets. Means over
seeds: PCA is deterministic, the Gaussian random projection has 5
seeds, and UMAP, t-SNE and PaCMAP 3 each, at library defaults.
\begin{center}
\small
\begin{tabular}{lccccc}
\hline
 & shuttle & wbc & ionosphere & fault & magic \\
\hline
PCA & 0.57 ($+0.34$) & 0.00 ($+0.91$) & 1.00 ($-0.05$) & 0.58 ($+0.37$) & 0.39 ($+0.56$) \\
Random & 0.40 ($+0.49$) & 0.37 ($+0.58$) & 0.66 ($+0.31$) & 0.59 ($+0.36$) & 0.38 ($+0.58$) \\
UMAP & 0.91 ($+0.02$) & 0.61 ($+0.35$) & 1.00 ($-0.05$) & 0.82 ($+0.11$) & 0.85 ($+0.10$) \\
t-SNE & 0.73 ($+0.18$) & 0.81 ($+0.15$) & 0.85 ($+0.10$) & 0.70 ($+0.25$) & 0.93 ($+0.02$) \\
PaCMAP & 0.85 ($+0.07$) & 0.97 ($-0.03$) & 1.00 ($-0.05$) & 0.98 ($-0.03$) & 0.89 ($+0.06$) \\
\hline
\end{tabular}
\end{center}

\begin{itemize}
  \item UMAP, t-SNE and PaCMAP hide 61 to 100\% of the top 5\%, more than PCA
    or a random projection on four of the five datasets. PaCMAP's shortlist
    is close to chance on all five (AO at most 0.09).
  \item PCA hides from 0\% (wbc) to 100\% (ionosphere). The failures fall as
    PCA keeps more dimensions and vanish by full dimension on all five.
  \item A random projection hides less than PCA on shuttle and ionosphere,
    because PCA discards low-variance directions whatever they hold.
  \item Aggregate correlations miss the shortlist. On ionosphere PCA keeps a
    Spearman correlation of 0.56, and 0.55 for the weighted correlation of
    Schubert et al.\ (2012), while its 2D shortlist shares nothing with the
    original one.
  \item Shuttle is not a control for PCA, which hides 57\% of its top 5\%:
    its median LID is about 2 locally, but it spans six linear dimensions.
    It is also lattice data: all nine features are integers, and its 2D
    pictures under PCA, UMAP, t-SNE and PaCMAP each break into
    1,188 to 1,511 small islands of near-identical points
    (169 under a random projection), against 5 to 65 for magic.
  \item The neighbour-based methods gather isolated outliers into tight
    clumps. On wbc and fault their verdict depends on the neighbourhood size
    used to read the picture; on shuttle, ionosphere and magic it does not.
\end{itemize}

\subsection*{Open, for the supervisor}

\begin{enumerate}
  \item A resampling null per point, with permutation kept only for the
    chance level of AO, because the scDEED-style permutation null calls
    almost nothing. Is that acceptable?
  \item Budget: the top 5\% of points, which needs no labels, instead of the
    number of labelled anomalies?
  \item How to read the projection: the same $k$ as the original space (the
    proposal's rule), a profile over several $k$, or a multi-scale score?
  \item Keep shuttle as the case where local and global dimension differ?
  \item Lattice data: keep shuttle and wbc as they are and report the effect?
    The discrete fields of network data will raise the same question.
  \item Ranks, as now, or normalised scores (Kriegel et al.\ 2011) for
    version 2?
\end{enumerate}
 in main.tex.
% Numbers: results/week4_measure.json and results/week4_lattice.json,
% written by src/explore_measure.py and src/lattice_check.py.

\section*{Week 4: outlier faithfulness, version 1}

Built and tested on the exploration set only: shuttle, wbc, ionosphere,
fault and magic. The repository calls these five the development set. No
benchmark dataset was used. Version 1 will change in week 5 and is frozen
in week 12.

\subsection*{Definition}

Let $x_1, \dots, x_n \in \mathbb{R}^D$ be the data and
$y_1, \dots, y_n \in \mathbb{R}^d$ a projection of it. The same detector
scores both spaces: $a_i$ is the distance from $x_i$ to its $k$-th nearest
neighbour among the other points, and $b_i$ the same for $y_i$, with
$k = 20$. Outlying in the projection therefore means ranked high by the
detector run on the projected coordinates.

\paragraph{Ranks.} $\rho_a(i)$ is the descending rank of $a_i$, with
$\rho_a(i) = 1$ for the most outlying point. Scores within a relative
tolerance of $10^{-9}$ count as tied, and tied points share the average of
the positions they span. $\rho_b$ is defined the same way.

\paragraph{Displacement,} in bits:
\begin{equation*}
  \lambda_i = \log_2 \frac{\rho_a(i)}{\rho_b(i)} .
\end{equation*}
A negative value means the projection moved the point towards the crowd, a
positive value towards the edge. $|\lambda_i|$ counts rank doublings, so a
move from rank 1 to rank 64 is $-6$ and a move from rank 1001 to rank 1064
is $-0.09$. The top of the ranking, where an analyst looks, dominates.

\paragraph{Shortlists.} With budget $m = \lceil 0.05\, n \rceil$,
$S_a = \{ i : \rho_a(i) \le m \}$ and $S_b = \{ i : \rho_b(i) \le m \}$. A
tied group shares one rank, so it is wholly in or wholly out.

\paragraph{Noise null.} For $r = 1, \dots, 100$, draw a random half $H_r$
of the points and score every point by its distance to the $(k/2)$-th
nearest neighbour within $H_r$, itself excluded. Ranking all $n$ points
gives $\rho^{(r)}$. The null for point $i$ is the set of differences
$D_i^{(rs)} = \log_2 \rho^{(r)}(i) - \log_2 \rho^{(s)}(i)$ over all pairs
$r < s$, and
\begin{equation*}
  p_i = \tfrac{1}{2} \Pr_{r<s}\bigl( |D_i^{(rs)}| \ge |\lambda_i| \bigr)
\end{equation*}
is the one-sided p-value in the direction the point moved. The null asks how
far the point's rank moves between two equally valid measurements of the
original data. It does not involve the projection, so it is computed once
per dataset.

\paragraph{Calls,} at level $\alpha = 0.05$:
\begin{center}
\begin{tabular}{lll}
\hline
category & rule & meaning \\
\hline
retained & $i \in S_a \cap S_b$ & kept on the shortlist \\
masked & $i \in S_a \setminus S_b$ and $p_i < \alpha$ & hidden, beyond noise \\
phantom & $i \in S_b \setminus S_a$ and $p_i < \alpha$ & promoted, beyond noise \\
ambiguous & crossed, $p_i \ge \alpha$ & crossed within noise \\
\hline
\end{tabular}
\end{center}

\paragraph{Aggregates.}
\begin{equation*}
  \mathrm{MR} = \frac{|M|}{|S_a|}, \qquad
  \mathrm{PR} = \frac{|P|}{|S_b|}, \qquad
  \mathrm{AO} = \frac{O - E}{\min(|S_a|, |S_b|) - E},
\end{equation*}
where $M$ and $P$ are the masked and phantom points, $O = |S_a \cap S_b|$,
and $E = |S_a|\,|S_b| / n$ is the overlap expected when the projected scores
are permuted across points. MR is the share of the detector's shortlist that
the projection hides, PR the share of the projection's shortlist that is not
outlying in the data. AO is 1 for identical shortlists and 0 at chance, the
same adjustment Campos et al.\ (2016) apply to precision at $n$.

\paragraph{Labels} are not used by the measure. A labelled anomaly that is
not outlying in the original space cannot be masked, which is the
conditioning on original-space outlyingness the proposal asks for. Labels
only say what a failure costs: a masked anomaly is a detection lost to the
projection.

\subsection*{Checks with known answers}

All 44 checks in \texttt{src/validate.py} pass. Among them: an identity
projection and a rotation move nothing; swapping the two spaces negates
$\lambda$ and swaps masked with phantom; a monotone rescoring changes
nothing; shuffled projected scores give an AO of 0 on average; comparing two
fresh re-measurements of the original data calls 2 to 3\% of crossings at
$\alpha = 5\%$; and of 20 outliers planted in the directions PCA discards,
17 to 20 are called masked over five seeds, while all 20 planted in the
plane PCA keeps are retained.

\subsection*{Results in two dimensions, top 5\%}

Masking rate MR, with the adjusted overlap AO in brackets. Means over
seeds: PCA is deterministic, the Gaussian random projection has 5
seeds, and UMAP, t-SNE and PaCMAP 3 each, at library defaults.
\begin{center}
\small
\begin{tabular}{lccccc}
\hline
 & shuttle & wbc & ionosphere & fault & magic \\
\hline
PCA & 0.57 ($+0.34$) & 0.00 ($+0.91$) & 1.00 ($-0.05$) & 0.58 ($+0.37$) & 0.39 ($+0.56$) \\
Random & 0.40 ($+0.49$) & 0.37 ($+0.58$) & 0.66 ($+0.31$) & 0.59 ($+0.36$) & 0.38 ($+0.58$) \\
UMAP & 0.91 ($+0.02$) & 0.61 ($+0.35$) & 1.00 ($-0.05$) & 0.82 ($+0.11$) & 0.85 ($+0.10$) \\
t-SNE & 0.73 ($+0.18$) & 0.81 ($+0.15$) & 0.85 ($+0.10$) & 0.70 ($+0.25$) & 0.93 ($+0.02$) \\
PaCMAP & 0.85 ($+0.07$) & 0.97 ($-0.03$) & 1.00 ($-0.05$) & 0.98 ($-0.03$) & 0.89 ($+0.06$) \\
\hline
\end{tabular}
\end{center}

\begin{itemize}
  \item UMAP, t-SNE and PaCMAP hide 61 to 100\% of the top 5\%, more than PCA
    or a random projection on four of the five datasets. PaCMAP's shortlist
    is close to chance on all five (AO at most 0.09).
  \item PCA hides from 0\% (wbc) to 100\% (ionosphere). The failures fall as
    PCA keeps more dimensions and vanish by full dimension on all five.
  \item A random projection hides less than PCA on shuttle and ionosphere,
    because PCA discards low-variance directions whatever they hold.
  \item Aggregate correlations miss the shortlist. On ionosphere PCA keeps a
    Spearman correlation of 0.56, and 0.55 for the weighted correlation of
    Schubert et al.\ (2012), while its 2D shortlist shares nothing with the
    original one.
  \item Shuttle is not a control for PCA, which hides 57\% of its top 5\%:
    its median LID is about 2 locally, but it spans six linear dimensions.
    It is also lattice data: all nine features are integers, and its 2D
    pictures under PCA, UMAP, t-SNE and PaCMAP each break into
    1,188 to 1,511 small islands of near-identical points
    (169 under a random projection), against 5 to 65 for magic.
  \item The neighbour-based methods gather isolated outliers into tight
    clumps. On wbc and fault their verdict depends on the neighbourhood size
    used to read the picture; on shuttle, ionosphere and magic it does not.
\end{itemize}

\subsection*{Open, for the supervisor}

\begin{enumerate}
  \item A resampling null per point, with permutation kept only for the
    chance level of AO, because the scDEED-style permutation null calls
    almost nothing. Is that acceptable?
  \item Budget: the top 5\% of points, which needs no labels, instead of the
    number of labelled anomalies?
  \item How to read the projection: the same $k$ as the original space (the
    proposal's rule), a profile over several $k$, or a multi-scale score?
  \item Keep shuttle as the case where local and global dimension differ?
  \item Lattice data: keep shuttle and wbc as they are and report the effect?
    The discrete fields of network data will raise the same question.
  \item Ranks, as now, or normalised scores (Kriegel et al.\ 2011) for
    version 2?
\end{enumerate}
 in main.tex.
% Numbers: results/week4_measure.json and results/week4_lattice.json,
% written by src/explore_measure.py and src/lattice_check.py.

\section*{Week 4: outlier faithfulness, version 1}

Built and tested on the exploration set only: shuttle, wbc, ionosphere,
fault and magic. The repository calls these five the development set. No
benchmark dataset was used. Version 1 will change in week 5 and is frozen
in week 12.

\subsection*{Definition}

Let $x_1, \dots, x_n \in \mathbb{R}^D$ be the data and
$y_1, \dots, y_n \in \mathbb{R}^d$ a projection of it. The same detector
scores both spaces: $a_i$ is the distance from $x_i$ to its $k$-th nearest
neighbour among the other points, and $b_i$ the same for $y_i$, with
$k = 20$. Outlying in the projection therefore means ranked high by the
detector run on the projected coordinates.

\paragraph{Ranks.} $\rho_a(i)$ is the descending rank of $a_i$, with
$\rho_a(i) = 1$ for the most outlying point. Scores within a relative
tolerance of $10^{-9}$ count as tied, and tied points share the average of
the positions they span. $\rho_b$ is defined the same way.

\paragraph{Displacement,} in bits:
\begin{equation*}
  \lambda_i = \log_2 \frac{\rho_a(i)}{\rho_b(i)} .
\end{equation*}
A negative value means the projection moved the point towards the crowd, a
positive value towards the edge. $|\lambda_i|$ counts rank doublings, so a
move from rank 1 to rank 64 is $-6$ and a move from rank 1001 to rank 1064
is $-0.09$. The top of the ranking, where an analyst looks, dominates.

\paragraph{Shortlists.} With budget $m = \lceil 0.05\, n \rceil$,
$S_a = \{ i : \rho_a(i) \le m \}$ and $S_b = \{ i : \rho_b(i) \le m \}$. A
tied group shares one rank, so it is wholly in or wholly out.

\paragraph{Noise null.} For $r = 1, \dots, 100$, draw a random half $H_r$
of the points and score every point by its distance to the $(k/2)$-th
nearest neighbour within $H_r$, itself excluded. Ranking all $n$ points
gives $\rho^{(r)}$. The null for point $i$ is the set of differences
$D_i^{(rs)} = \log_2 \rho^{(r)}(i) - \log_2 \rho^{(s)}(i)$ over all pairs
$r < s$, and
\begin{equation*}
  p_i = \tfrac{1}{2} \Pr_{r<s}\bigl( |D_i^{(rs)}| \ge |\lambda_i| \bigr)
\end{equation*}
is the one-sided p-value in the direction the point moved. The null asks how
far the point's rank moves between two equally valid measurements of the
original data. It does not involve the projection, so it is computed once
per dataset.

\paragraph{Calls,} at level $\alpha = 0.05$:
\begin{center}
\begin{tabular}{lll}
\hline
category & rule & meaning \\
\hline
retained & $i \in S_a \cap S_b$ & kept on the shortlist \\
masked & $i \in S_a \setminus S_b$ and $p_i < \alpha$ & hidden, beyond noise \\
phantom & $i \in S_b \setminus S_a$ and $p_i < \alpha$ & promoted, beyond noise \\
ambiguous & crossed, $p_i \ge \alpha$ & crossed within noise \\
\hline
\end{tabular}
\end{center}

\paragraph{Aggregates.}
\begin{equation*}
  \mathrm{MR} = \frac{|M|}{|S_a|}, \qquad
  \mathrm{PR} = \frac{|P|}{|S_b|}, \qquad
  \mathrm{AO} = \frac{O - E}{\min(|S_a|, |S_b|) - E},
\end{equation*}
where $M$ and $P$ are the masked and phantom points, $O = |S_a \cap S_b|$,
and $E = |S_a|\,|S_b| / n$ is the overlap expected when the projected scores
are permuted across points. MR is the share of the detector's shortlist that
the projection hides, PR the share of the projection's shortlist that is not
outlying in the data. AO is 1 for identical shortlists and 0 at chance, the
same adjustment Campos et al.\ (2016) apply to precision at $n$.

\paragraph{Labels} are not used by the measure. A labelled anomaly that is
not outlying in the original space cannot be masked, which is the
conditioning on original-space outlyingness the proposal asks for. Labels
only say what a failure costs: a masked anomaly is a detection lost to the
projection.

\subsection*{Checks with known answers}

All 44 checks in \texttt{src/validate.py} pass. Among them: an identity
projection and a rotation move nothing; swapping the two spaces negates
$\lambda$ and swaps masked with phantom; a monotone rescoring changes
nothing; shuffled projected scores give an AO of 0 on average; comparing two
fresh re-measurements of the original data calls 2 to 3\% of crossings at
$\alpha = 5\%$; and of 20 outliers planted in the directions PCA discards,
17 to 20 are called masked over five seeds, while all 20 planted in the
plane PCA keeps are retained.

\subsection*{Results in two dimensions, top 5\%}

Masking rate MR, with the adjusted overlap AO in brackets. Means over
seeds: PCA is deterministic, the Gaussian random projection has 5
seeds, and UMAP, t-SNE and PaCMAP 3 each, at library defaults.
\begin{center}
\small
\begin{tabular}{lccccc}
\hline
 & shuttle & wbc & ionosphere & fault & magic \\
\hline
PCA & 0.57 ($+0.34$) & 0.00 ($+0.91$) & 1.00 ($-0.05$) & 0.58 ($+0.37$) & 0.39 ($+0.56$) \\
Random & 0.40 ($+0.49$) & 0.37 ($+0.58$) & 0.66 ($+0.31$) & 0.59 ($+0.36$) & 0.38 ($+0.58$) \\
UMAP & 0.91 ($+0.02$) & 0.61 ($+0.35$) & 1.00 ($-0.05$) & 0.82 ($+0.11$) & 0.85 ($+0.10$) \\
t-SNE & 0.73 ($+0.18$) & 0.81 ($+0.15$) & 0.85 ($+0.10$) & 0.70 ($+0.25$) & 0.93 ($+0.02$) \\
PaCMAP & 0.85 ($+0.07$) & 0.97 ($-0.03$) & 1.00 ($-0.05$) & 0.98 ($-0.03$) & 0.89 ($+0.06$) \\
\hline
\end{tabular}
\end{center}

\begin{itemize}
  \item UMAP, t-SNE and PaCMAP hide 61 to 100\% of the top 5\%, more than PCA
    or a random projection on four of the five datasets. PaCMAP's shortlist
    is close to chance on all five (AO at most 0.09).
  \item PCA hides from 0\% (wbc) to 100\% (ionosphere). The failures fall as
    PCA keeps more dimensions and vanish by full dimension on all five.
  \item A random projection hides less than PCA on shuttle and ionosphere,
    because PCA discards low-variance directions whatever they hold.
  \item Aggregate correlations miss the shortlist. On ionosphere PCA keeps a
    Spearman correlation of 0.56, and 0.55 for the weighted correlation of
    Schubert et al.\ (2012), while its 2D shortlist shares nothing with the
    original one.
  \item Shuttle is not a control for PCA, which hides 57\% of its top 5\%:
    its median LID is about 2 locally, but it spans six linear dimensions.
    It is also lattice data: all nine features are integers, and its 2D
    pictures under PCA, UMAP, t-SNE and PaCMAP each break into
    1,188 to 1,511 small islands of near-identical points
    (169 under a random projection), against 5 to 65 for magic.
  \item The neighbour-based methods gather isolated outliers into tight
    clumps. On wbc and fault their verdict depends on the neighbourhood size
    used to read the picture; on shuttle, ionosphere and magic it does not.
\end{itemize}

\subsection*{Open, for the supervisor}

\begin{enumerate}
  \item A resampling null per point, with permutation kept only for the
    chance level of AO, because the scDEED-style permutation null calls
    almost nothing. Is that acceptable?
  \item Budget: the top 5\% of points, which needs no labels, instead of the
    number of labelled anomalies?
  \item How to read the projection: the same $k$ as the original space (the
    proposal's rule), a profile over several $k$, or a multi-scale score?
  \item Keep shuttle as the case where local and global dimension differ?
  \item Lattice data: keep shuttle and wbc as they are and report the effect?
    The discrete fields of network data will raise the same question.
  \item Ranks, as now, or normalised scores (Kriegel et al.\ 2011) for
    version 2?
\end{enumerate}
